\documentclass[10pt,twocolumn]{article}
\usepackage[T1]{fontenc}
\usepackage[utf8]{inputenc}
\usepackage{lmodern}
\usepackage[margin=1.8cm,top=2.4cm,bottom=2cm,columnsep=0.6cm]{geometry}
\usepackage{fancyhdr}
\usepackage{titlesec}
\usepackage{booktabs}
\usepackage{amsmath,amssymb,amsfonts}
\usepackage{microtype}
\usepackage{xcolor}
\usepackage{mdframed}
\usepackage{parskip}
\usepackage{enumitem}
\usepackage{hyperref}

\definecolor{rulered}{RGB}{180,30,30}
\definecolor{boxgray}{RGB}{245,245,245}
\definecolor{keywordgray}{RGB}{100,100,100}

\pagestyle{fancy}
\fancyhf{}
\fancyhead[L]{\textsc{\small Claude Mag}}
\fancyhead[C]{\small\textit{Machine Learning Research Digest}}
\fancyhead[R]{\small 2026-08-03}
\fancyfoot[C]{\small\thepage}
\renewcommand{\headrulewidth}{0.8pt}

\titleformat{\section}{\normalsize\bfseries\color{rulered}}{}{0pt}{}%
  [\vspace{-4pt}\color{rulered}\rule{\columnwidth}{0.4pt}]
\titlespacing{\section}{0pt}{8pt}{4pt}

\hypersetup{colorlinks=true,urlcolor=rulered,linkcolor=black}

\begin{document}
\setlength{\parindent}{0pt}
\setlength{\parskip}{4pt}

{\small\color{keywordgray}\textbf{Keywords:}
  vision-language-action \textbullet{}
  robot policy \textbullet{}
  real-time inference \textbullet{}
  efficient VLA}\par\vspace{4pt}

{\LARGE\bfseries\raggedright
  TurboVLA: Real-Time Vision-Language-Action Model
  at 32\,Hz on an RTX\,4090 with $<$1\,GB VRAM}\par\vspace{4pt}

{\large\itshape\raggedright
  Eliminating the LLM Bottleneck in Robot Policies via Direct
  V$+$L$\to$A Mapping Yields 97.7\% on LIBERO with Only
  0.2B Parameters}\par\vspace{6pt}

{\small\textbf{By} Hengyi Xie, Chenfei Yao, Xianjin Wu, Xuanyang Xi,
  Yiping Tang, Di Xu, Yingying Zhu, Dingkang Liang, Xiang Bai, Han Ding
  (Huazhong Univ.\ of Science and Technology; Huawei Technologies)}\par\vspace{2pt}

{\small\color{keywordgray}arXiv:2607.27205 \textbullet{} July 30, 2026}
\par\vspace{2pt}\noindent{\color{rulered}\rule{\columnwidth}{0.6pt}}\vspace{6pt}

State-of-the-art robot manipulation models route all perception through a
billion-parameter language model---incurring hundreds of milliseconds of
latency and gigabytes of VRAM at every policy step. \textbf{TurboVLA}
breaks from this convention, replacing the V$\to$L$\to$A pathway with a
direct V$+$L$\to$A mapping using a compact bidirectional interaction module.
The result runs at 32\,Hz on a consumer RTX\,4090 with $<$1\,GB VRAM while
achieving 97.7\% average success on LIBERO---matching models 35$\times$
larger.

\begin{mdframed}[backgroundcolor=boxgray,linecolor=rulered,linewidth=1pt,
    innerleftmargin=6pt,innerrightmargin=6pt,
    innertopmargin=6pt,innerbottommargin=6pt]
\textbf{\small AT A GLANCE}\par\vspace{4pt}
\begin{description}[leftmargin=0pt,labelsep=4pt,itemsep=2pt,parsep=0pt,
    font=\small\bfseries]
  \item[Problem:] LLM-centric VLA policies incur prohibitive compute and
    memory overhead at every inference step, blocking real-time deployment.
  \item[Why it matters:] Real-time control at 32\,Hz on consumer hardware
    is a prerequisite for affordable robot deployment outside research labs.
  \item[Method:] Independent visual and language encoders fused via
    lightweight bidirectional cross-attention; continuous action chunks
    predicted by a compact decoder.
  \item[Result:] 97.7\% on LIBERO, 31.2\,ms per step, 0.9\,GB VRAM,
    0.2B parameters.
  \item[Evidence:] Matches or exceeds OpenVLA-7B and Pi-0 (3.3B) while
    running 7--35$\times$ faster and consuming 7--16$\times$ less memory.
\end{description}
\end{mdframed}

\section{The Big Picture}

Vision-language-action models represent the current frontier of generalist
robot policies: by building on pre-trained vision-language models, they
inherit rich semantic representations that generalize across diverse
manipulation tasks. Papers like OpenVLA and Pi-0 have demonstrated that
leveraging multi-billion-parameter LLM backbones as the core of a robot
policy can yield impressive task success rates.

The bottleneck, however, is clear: inference in these models runs at
2--10\,Hz on research-grade GPU hardware, requiring tens of gigabytes of
VRAM. Consumer robots, surgical systems, and industrial manipulators
demand 30--100\,Hz control with modest compute budgets. TurboVLA
challenges the architectural assumption that the LLM must be in the
critical path of every policy inference.

\section{Why Existing Approaches Fall Short}

The dominant V$\to$L$\to$A paradigm processes visual observations by
projecting them into the LLM's token space: visual patches become
``visual tokens'' that attend to language tokens in the LLM's self-attention
layers. This design has three costs:

\textbf{Quadratic attention complexity.} The joint visual-language sequence
has $N_v + N_l$ tokens; self-attention costs $O((N_v + N_l)^2)$ per layer.
For $N_v = 256$ visual tokens and $N_l = 30$ language tokens, this is
$\sim 80{,}000$ token pairs per layer, repeated across all 32+ LLM layers.

\textbf{Memory dominated by LLM weights.} Even a 7B-parameter LLM at
4-bit quantization occupies $\sim 3.5$\,GB---more than most embedded
compute platforms support.

\textbf{No separation of concerns.} The LLM processes both semantic
instruction understanding and low-level observation encoding in the same
computational graph, making it impossible to specialize or prune either
component independently.

\section{Core Mechanism}

TurboVLA introduces three specialized modules that replace the monolithic
LLM backbone:

\begin{enumerate}[itemsep=1pt,parsep=0pt]
  \item \textbf{Visual encoder (ViT-S/16):} Encodes the observation image
    into $N_v = 196$ patch tokens $V \in \mathbb{R}^{N_v \times d_v}$.
  \item \textbf{Language encoder (CLIP text, frozen):} Encodes the
    instruction into $N_l \approx 30$ tokens $L \in \mathbb{R}^{N_l \times d_l}$.
  \item \textbf{Bidirectional Vision-Language Interaction (BiVLI):}
    Applies two rounds of alternating cross-attention, letting $V$ attend
    to $L$ and $L$ attend to $V$, producing aligned representations $V'$
    and $L'$.
\end{enumerate}

The BiVLI module contains only $\sim$50M parameters. The action decoder
takes $[\mathrm{pool}(V') \oplus \mathrm{pool}(L')]$ and predicts a
chunk of $H = 10$ continuous actions.

\section{Technical Formulation}

For round $r \in \{1, 2\}$ of BiVLI:
\[
  V^{(r)} = V^{(r-1)} + \mathrm{MHA}_{V\leftarrow L}(Q{=}V^{(r-1)},\; K{=}L^{(r-1)},\; V{=}L^{(r-1)})
\]
\[
  L^{(r)} = L^{(r-1)} + \mathrm{MHA}_{L\leftarrow V}(Q{=}L^{(r-1)},\; K{=}V^{(r)},\; V{=}V^{(r)})
\]
where $\mathrm{MHA}$ is multi-head attention. The interaction cost is
$O(N_v \cdot N_l)$ per round---$30\times$ lower than full self-attention
over the combined sequence.

The combined representation $z = \mathrm{pool}(V^{(2)}) \oplus \mathrm{pool}(L^{(2)})$
is decoded to an action chunk:
\[
  \hat{a}_{1:H} = \mathrm{Decoder}(z), \qquad
  \hat{a}_{1:H} \in \mathbb{R}^{H \times d_a}
\]

Training minimizes the behavior cloning loss:
\[
  \mathcal{L}_{\text{BC}} = \sum_{h=1}^{H} \bigl\|\hat{a}_h - a_h^*\bigr\|_2^2
\]
over expert demonstrations.

\section{Learning or Inference Procedure}

Training uses two stages. In the first (pretraining) stage, the visual
encoder is initialized from CLIP ViT-S/16 and the BiVLI module is
randomly initialized; both are trained on the OpenX-Embodiment dataset
($\sim$1M demonstrations). In the second (fine-tuning) stage, the full
model including the visual encoder is fine-tuned on task-specific
demonstrations in the target environment (LIBERO, $\sim$500
demonstrations per suite). The CLIP text encoder is frozen throughout.

At inference, the visual encoder, BiVLI, and action decoder run in a
single forward pass per time step---no autoregressive decoding, no KV
cache management, no LLM overhead.

\section{What the Guarantee Says}

By design, the backbone parameters $p_\theta$ are never called at
inference time. The only per-step compute is ViT-S forward (12\,ms),
BiVLI two-round cross-attention (11\,ms), and MLP action decoder (8\,ms),
totaling 31.2\,ms---safely below the 33\,ms deadline for 30\,Hz control.
This is a structural bound, not a statistical approximation: no matter
how complex the manipulation task, the per-step latency is fixed by
architecture rather than by task difficulty.

\section{Experimental Findings}

\textbf{LIBERO benchmark (5 suites, higher is better):}

\begin{center}\small
\begin{tabular}{lcccc}
\toprule
Model & Params & Avg. & ms/step & VRAM \\
\midrule
OpenVLA-7B & 7B & 96.1\% & 218 & 14.0\,GB \\
Pi-0 (distil.) & 3.3B & 97.2\% & 82 & 6.4\,GB \\
TurboVLA & \textbf{0.2B} & \textbf{97.7\%} & \textbf{31} & \textbf{0.9\,GB} \\
\bottomrule
\end{tabular}
\end{center}

TurboVLA achieves the best average success rate while using 35$\times$
fewer parameters, running 7$\times$ faster, and consuming 7--16$\times$
less memory.

\textbf{Cross-embodiment:} Evaluated on Franka, xArm, and UR5, TurboVLA
shows less catastrophic forgetting between embodiments than OpenVLA,
attributed to the frozen CLIP text encoder providing a stable semantic
anchor during multi-embodiment fine-tuning.

\section{Ablations and Interpretation}

\textbf{BiVLI rounds:} $R=1$ round drops average success by 3.2\%; $R=4$
rounds adds only 0.4\% while doubling BiVLI latency. Two rounds is the
Pareto-optimal choice for accuracy/latency trade-off.

\textbf{Visual encoder size:} ViT-B/16 (3$\times$ more parameters) improves
success by 1.1\% but doubles VRAM. ViT-S/16 is appropriate for
resource-constrained deployment.

\textbf{Action chunk horizon:} $H=10$ gives the best task completion;
shorter horizons ($H=4$) hurt due to high-frequency control jitter and
longer horizons ($H=20$) hurt due to compounding prediction error.

\section*{Reference}

Hengyi Xie, Chenfei Yao, Xianjin Wu, Xuanyang Xi, Yiping Tang, Di Xu,
Yingying Zhu, Dingkang Liang, Xiang Bai, Han Ding.
\textbf{TurboVLA: Real-Time Vision-Language-Action Model at 32\,Hz on
an RTX\,4090 with $<$1\,GB VRAM}.
arXiv:2607.27205, July 30, 2026.
\url{https://arxiv.org/abs/2607.27205}

\end{document}
